\subsection{COMPOSITION OF ASYMPTOTIC EXPANSIONS}

Let f be a function in \(\mathcal{H}(\mathfrak{F},\mathbf{R})\) (resp. \(\mathcal{H}(\mathfrak{F},\mathbf{C})\) ), admitting an asymptotic expansion to precision \(g_{\alpha}\) with respect to a scale E, and having limit 0 along the filter with base F. On the other hand let h be a function with values in a normed space V over R (resp. C), defined on a neighbourhood of the point 0 in R (resp. C), and n times differentiable on this neighbourhood; then

\[
\mathbf {h} (t) = \mathbf {c} _ {0} + \mathbf {c} _ {1} t + \dots + \mathbf {c} _ {n} t ^ {n} + o (t ^ {n})
\]

on this neighbourhood (I, p.~21), whence, on a suitable set in \(\mathfrak{F}\)

\[
\mathbf {h} \circ f = \mathbf {c} _ {0} + \mathbf {c} _ {1} f + \dots + \mathbf {c} _ {n} f ^ {n} + o (f ^ {n}).
\]

We have seen, in n° 3, how to form an asymptotic expansion for \(c_{0} + c_{1}f + \cdots + c_{n}f^{n}\) to a precision \(g_{\rho}\) determined by the precision of the expansion of f; moreover, suppose that the coefficients of the asymptotic expansion of f are not all zero, and that \(a_{\gamma}g_{\gamma}\) is the principal part of f, and let \(g_{\sigma} = g_{\gamma}^{n}\) ; if \(\sigma < \rho\) one will have an expansion of \(h \circ f\) to precision \(g_{\sigma}\) on limiting the expansion of \(\sum_{k=0}^{n} c_{k}f^{k}\) to this precision; if, on the contrary, \(\rho \leqslant \sigma\) , then the expansion of \(\sum_{k=0}^{n} c_{k}f^{k}\) is also an expansion of \(h \circ f\) to precision \(g_{\rho}\) .

If all the terms of the asymptotic expansion of \(f\) are zero, and if \(g_{\alpha} \ll 1\) , then \(f \ll g_{\alpha}\) and so \(f^k \ll g_{\alpha}^k \ll g_{\alpha}\) for every integer \(k > 0\) ; if \(\mathbf{c}_m\) is the first coefficient of index \(> 0\) which is not zero (assuming that the \(\mathbf{c}_k\) for indices \(k > 0\) are not all 0), then \(\mathbf{c}_0\) is an asymptotic expansion of \(\mathbf{h} \circ f\) to precision \(g_{\alpha}^{m}\) .

In the remainder of this section we shall restrict ourselves to the case where the functions in \(\mathcal{E}\) have real values and are strictly positive on a set in \(\mathfrak{F}\) , and we shall consider asymptotic expansions only of functions in \(\mathcal{H}(\mathfrak{F},\mathbf{R})\) . Suppose first that for every function \(g\in \mathcal{E}\) and every real number \(\nu\) , \(g^{\nu}\) again belongs to \(\mathcal{E}\) : this condition is fulfilled, for example, by the scale of the \(x^{\alpha}\) or by that of the \(x^{\alpha}|\log x|^{\beta}\) ( \(\alpha\) and \(\beta\) being arbitrary real numbers) on a neighbourhood of \(+\infty\) or a neighbourhood of 0 in \(\mathbf{R}\) . This property implies that the quotient of two functions in \(\mathcal{E}\) again belongs to \(\mathcal{E}\) . This being so, from an asymptotic expansion relative to \(\mathcal{E}\) of a function \(f\in \mathcal{H}(\mathfrak{F},\mathbf{R})\) , to precision \(g_{\alpha}\) , one can derive an expansion for \(|f|^{\nu}\) for every real number \(\nu\) . Let us restrict ourselves to the case where the coefficients of the expansion of \(f\) are not all zero, and let \(a_{\gamma}g_{\gamma}\) be the principal part of \(f\) ; one can write \(|f|^{\nu} = |a_{\gamma}|^{\nu}g_{\gamma}^{\nu}(1 + h)^{\nu}\) , with

\[
h = \sum_ {\gamma <   \lambda \leqslant \alpha} \frac {a _ {\lambda}}{a _ {\gamma}} \frac {g _ {\lambda}}{g _ {\gamma}} + o \left(\frac {g _ {\alpha}}{g _ {\gamma}}\right).
\]

Under our hypotheses \(\sum_{\gamma < \lambda \leqslant \alpha} \frac{a_{\lambda}}{a_{\gamma}} \frac{g_{\lambda}}{g_{\gamma}}\) is an asymptotic expansion of \(h\) , to precision \(g_{\alpha} / g_{\gamma}\) ; since \(h\) tends to 0 along \(\mathfrak{F}\) the method described above gives an asymptotic expansion of \((1 + h)^{\nu}\) , then an expansion of \(|f|^{\nu}\) on multiplying by \(\left|a_{\gamma}\right|^{\nu} g_{\gamma}^{\nu}\) .

With the same hypotheses on \(f\) one can write

\[
\log | f | = \log \left| a _ {\gamma} g _ {\gamma} \right| + \log (1 + h)
\]

and \(\log(1+h)\) can be expanded, as has been said above, the function \(\log(1+t)\) being indefinitely differentiable on a neighbourhood of 0; if, further, \(\log g_{\gamma}\) admits an asymptotic expansion with respect to E, or with respect to a scale \(E_{1} \supset E\) , one obtains an asymptotic expansion for \(\log|f|\) on adding the two asymptotic expansions.

Example. We have \((1+x)^{1/x}=\exp\left(\frac{1}{x}\log(1+x)\right)\) ; when x tends to \(+\infty\) we have \(\log(1+x)=\log x+\log\left(1+\frac{1}{x}\right)\) , whence the asymptotic expansion of \(\frac{1}{x}\log(1+x)\) with respect to the scale of the \(x^{\alpha}(\log x)^{\beta}\) :

\[
\frac {1}{x} \log (1 + x) = \frac {\log x}{x} + \frac {1}{x ^ {2}} - \frac {1}{2 x ^ {3}} + o _ {1} \left(\frac {1}{x ^ {3}}\right).
\]

From this expansion, and from the Taylor expansion

\[
e ^ {u} = 1 + u + \frac {u ^ {2}}{2} + \frac {u ^ {3}}{6} + o (u ^ {3})
\]

on a neighbourhood of u = 0, one deduces the asymptotic expansion

\[
\begin{array}{r l} (1 + x) ^ {1 / x} & = 1 + \frac {\log x}{x} + \frac {1}{2} \frac {(\log x) ^ {2}}{x ^ {2}} + \frac {1}{x ^ {2}} \\ & \quad + \frac {1}{6} \frac {(\log x) ^ {3}}{x ^ {3}} + \frac {\log x}{x ^ {3}} - \frac {1}{2 x ^ {3}} + o _ {2} \left(\frac {1}{x ^ {3}}\right) \end{array}
\]

with respect to the scale \(x^{\alpha}(\log x)^{\beta}\) by the methods explained above.

Keeping the same hypotheses and notation, the asymptotic expansion of \(e^{f}\) does not pose new problems except when \(f \gg 1\) ; one must then distinguish two cases, as \(g_{\alpha} \gg 1\) or \(g_{\alpha} \preccurlyeq 1\) . In the first case, giving an expansion for f does not allow one to obtain a principal part for \(e^{f}\) relative to E, because one does not know in general whether the remainder \(r_{\alpha}\) tends to 0, that is, whether \(e^{r_{\alpha}}\) tends to 1. On the contrary, if \(g_{\alpha} \preccurlyeq 1\) then \(r_{\alpha} \ll 1\) and so \(e^{f} \sim \exp\left(\sum_{\lambda \leqslant \alpha} a_{\lambda} g_{\lambda}\right)\) . One can make this result more precise: let \(a_{\gamma} g_{\gamma}\) be the principal part of f and let \(\delta\) be the index (such that \(\gamma < \delta \leqslant \alpha\) ) for which \(g_{\delta} = 1\) ; we put \(f_{1} = \sum_{\lambda \leqslant \delta} a_{\lambda} g_{\lambda}\) , \(f_{2} = \sum_{\delta < \lambda \leqslant \alpha} a_{\lambda} g_{\lambda} + r_{\alpha}\) ; we have \(f = f_{1} + f_{2}\) , so \(e^{f} = e^{f_{1}} e^{f_{2}}\) , and the general method explained at the start of this subsection allows one to form an asymptotic expansion of \(e^{f_{2}}\) (starting from the Taylor expansion of \(e^{t}\) at the point t = 0). One will then have an asymptotic expansion for \(e^{f}\) if \(e^{f_{1}} = \prod_{\lambda \leqslant \delta} \exp(a_{\lambda} g_{\lambda})\) belongs to E, or to a scale \(E_{1}\) containing E.

Example. We have \(x^{x^{1/x}} = \exp\left(\log x. \exp\left(\frac{1}{x} \log x\right)\right)\) ; when x tends to \(+\infty\) one has \(\log x \ll x\) , whence the asymptotic expansion of \(\log x. \exp\left(\frac{1}{x} \log x\right)\) with respect to the scale \(x^{\alpha}(\log x)^{\beta}\) :

\[
\log x. \exp \left(\frac {1}{x} \log x\right) = \log x + \frac {(\log x) ^ {2}}{x} + \frac {1}{2} \frac {(\log x) ^ {3}}{x ^ {2}} + o \left(\frac {(\log x) ^ {3}}{x ^ {2}}\right).
\]

All the terms of this expansion, starting from the second, tend to 0; from this expansion and from the Taylor expansion \(e^u = 1 + u + u^2 / 2 + o(u^2)\) on a neighbourhood of \(u = 0\) one deduces

\[
x ^ {x ^ {1 / x}} = x + (\log x) ^ {2} + \frac {1}{2} \frac {(\log x) ^ {4}}{x} + \frac {1}{2} \frac {(\log x) ^ {3}}{x} + o \left(\frac {(\log x) ^ {3}}{x}\right).
\]

